% page 414 of the facsimile (image p-413 of the scan)
% block 152.4 x 233.3 mm ; left margin 8.0 mm
\pgc{-12.700mm}{0.000mm}{
\l{\cel{3}406}
\l{}
\l{\sou{onagro}. I. (\sou{zool.)}\cel{2}Asno sovaja. L.\cel{1}\sou{asinus onager}. II.}
\l{\cel{3}(\sou{militarto}). +Asiejo-mashino,uzita da la Romani anti-}
\l{\cel{3}qua, qua esis quaza balisto per qua onu lansis aden la}
\l{\cel{3}urbo +asiejata projektili diversa. - DEFIS.}
\l{}
\l{\sou{onaniar}. (\sou{netrans.})\cel{2}Produktar la orgasmo venerala per manuo}
\l{\cel{3}vice per la koito normala. - DEFIRS.}
\l{}
\l{\sou{ondo.}\cel{3}Amaso de aquo qua superstaceskas e retrofalas sur}
\l{\cel{3}la surfaco di maro, lago, fluvio, agitata da vento. -}
\l{\cel{3}EFIS.}
\l{}
\l{\sou{onixo}.\cel{2}Agato delikata, kun strati interparalela diversa-}
\l{\cel{3}kolora. - DEFIRS.}
\l{}
\l{\sou{onklo}. Frato di onua patro. - DEF.}
\l{}
\l{\sou{onomatopeo}. (\sou{gram.)}\cel{3}Vorto formacita per imito-harmonio.}
\l{\cel{3}(ex.: gluglo, zumo). - DEFIS.}
\l{}
\l{\sou{-ont-}\cel{4}(\sou{gram.)}\cel{4}Dezinenco di la participo futura ak-}
\l{\cel{3}tiva.}
\l{}
\l{\sou{ontogenezo.}\cel{2}(\sou{biol.)}\cel{2}Developo di la individuo, opoze a}
\l{\cel{3}\textquotedbl{}filogenezo\textquotedbl{}, o developo di la speco. - DEFIS.}
\l{}
\l{\sou{ontologio}. I. (\sou{biol.)}\cel{3}Parto di biologio, qua traktas}
\l{\cel{3}la fenomeni individuala quin prizentas organismo dum}
\l{\cel{3}lua developo, lua adulteso e lua oldeso. - II. (\sou{filoz}.)}
\l{\cel{3}Cienco metafizikala qua retroiras til la esenco, la}
\l{\cel{3}kauzo prima di la ento. - DEFIRS.}
\l{}
\l{\sou{onyono}. (\sou{bot.}) Planto leguma, di qua la radiko bulba -}
\l{\cel{3}qua odoras e saporas forte - uzesas en la preparaji}
\l{\cel{3}nutrala. L.\cel{1}\sou{allkum cepa}.\cel{2}EF.}
\l{}
\l{\sou{oocito}.\cel{2}(\sou{biol.}) Ovulo nematura nek fekundigebla, qua}
\l{\cel{3}produktas (per du dupartigi) la ovulo matura e fekun-}
\l{\cel{3}digebla, e la globuli polala. - DEFIRS.}
\l{}
\l{\sou{oogenezo}.\cel{3}(\sou{biol.)}\cel{2}Formaco di la ovulo matura e fe-}
\l{\cel{3}kundigebla. - DEFIRS.}
\l{}
\l{\sou{oogonio}. (\sou{biol.)}\cel{3}Formaco di la celulo en qua aparas}
\l{\cel{3}la elementi feminala, o\cel{1}\sou{oosferi}, che multa vejetanti.}
\l{}
\l{\sou{oolito}. (\sou{geol.)}\cel{3}Kalkajo ek mikra grani ovoida qui kelke}
\l{\cel{3}similesas fish-ovi. - DEFIRS.}
\l{}
\l{\sou{oosfero}. (\sou{biol.)}\cel{2}Elemento feminala qua, pos fekundigesir}
\l{\cel{3}da la elemento maskulala (polinido od anterozoido)}
\l{\cel{3}genitas la ovo - departo-punto di nova ento vejetanta.}
}
